\chapter{Group Introduction}
\label{chap:group-introduction}

\section{Group Information}

This section lists the group name, the student IDs, and the full member list of the team that carried out this project.

\begin{table}[h!]
\centering
\caption{Group information}
\begin{tabular}{ll}
\toprule
Item & Detail \\
\midrule
Group name & \fillin{group name} \\
Class & \fillin{class code} \\
Course & Introduction to Artificial Intelligence \\
Instructor & \fillin{instructor name} \\
Number of members & \fillin{number} \\
\bottomrule
\end{tabular}
\end{table}

\begin{table}[h!]
\centering
\caption{Member list}
\begin{tabular}{clc}
\toprule
No. & Full Name & Student ID \\
\midrule
1 & \fillin{member 1 name} & \fillin{ID} \\
2 & \fillin{member 2 name} & \fillin{ID} \\
3 & \fillin{member 3 name} & \fillin{ID} \\
4 & \fillin{member 4 name} & \fillin{ID} \\
5 & \fillin{member 5 name} & \fillin{ID} \\
\bottomrule
\end{tabular}
\end{table}

\section{Specific Contributions of Each Member}

The table below records what each member was responsible for, so that the division of work is clear and traceable. Replace the placeholder text with the actual tasks each person completed, for example dataset collection, graph construction, a specific search algorithm, the GUI, or the report writing.

\begin{table}[h!]
\centering
\caption{Contribution of each member}
\begin{tabular}{p{2.8cm}p{7cm}p{2cm}}
\toprule
Member & Contribution & Completion \\
\midrule
\fillin{name 1} & \fillin{tasks performed, for example graph modeling and dataset preparation} & \fillin{percent} \\
\fillin{name 2} & \fillin{tasks performed, for example implementation of BFS and DFS} & \fillin{percent} \\
\fillin{name 3} & \fillin{tasks performed, for example implementation of UCS and A star} & \fillin{percent} \\
\fillin{name 4} & \fillin{tasks performed, for example GUI development} & \fillin{percent} \\
\fillin{name 5} & \fillin{tasks performed, for example report writing and testing} & \fillin{percent} \\
\bottomrule
\end{tabular}
\end{table}

\section{Overall Completion Level of Each Project Requirement}

This table gives an honest self assessment of how complete each part of the assignment is. Use it to flag anything that is only partially done, so the reader knows what to expect before reading the rest of the report.

\begin{table}[h!]
\centering
\caption{Completion status of each requirement}
\begin{tabular}{p{7.5cm}p{2cm}p{4cm}}
\toprule
Requirement & Status & Note \\
\midrule
Graph modeling of the traffic network & \fillin{done / partial} & \fillin{note} \\
Dataset collection and preprocessing & \fillin{done / partial} & \fillin{note} \\
Search algorithm implementation & \fillin{done / partial} & \fillin{note} \\
Heuristic design for informed search & \fillin{done / partial} & \fillin{note} \\
GUI development & \fillin{done / partial} & \fillin{note} \\
Multi location route optimization & \fillin{done / partial} & \fillin{note} \\
Testing and performance comparison & \fillin{done / partial} & \fillin{note} \\
Report writing & \fillin{done / partial} & \fillin{note} \\
\bottomrule
\end{tabular}
\end{table}
